\section{Cut-cone evaluation, random-order capacity, and negative-bracket cut examples}\label{sec:sweep}

\subsection{Cut-causal cone evaluation on frozen \texttt{Transpose.TM} fibres}\label{sec:sweep-trees}

Let $M$ be a deterministic multitape Turing machine (\leanname{Transpose.TM}) with $1 \le M.k$
(available without loss of generality for any transposer by \Cref{lem:pad-dummy-tape}), and let
$C \subseteq \Sigma_d^{d\times d}$ be a frozen-trace $\Gamma$-cut fibre
(\lean{IsFrozenTraceCutFibre}) over a step horizon $T$: all $X \in C$ halt at step $T$ with output
$X^{\T}$, follow a common input/work head trajectory $(\mathrm{inPos}, \mathrm{workPos})$ and output
emission schedule $\mathrm{emitStep}$, and agree on the post-step control state and written
work-tape symbols at every cut step $u \in \Gamma \subseteq \{0,\dots,T-1\}$. Along
$\mathrm{workPos}$, a direct causal step $u \to v$ (\lean{TMDepStep}) holds for $u < v < T$ when
either $v = u + 1$ (the control-state spine edge) or $u$ is the most recent step visiting the cell
scanned by one of the $M.k$ work heads at step $v$ (\lean{IsLastWorkVisit}). A descending path is
$\Gamma$-free (\lean{IsCutFreeCausalPath}) if none of its strict predecessors lies in $\Gamma$; the
resulting $\Gamma$-truncated causal cone of step $v$ is \lean{trueConeSteps} (\lean{InCutCausalCone}),
and its syntactic input-cell support is \lean{cutConeInputCells}.

The causal-cone determination principle and out-degree cone-size bound go back in text to
\cite[Lemma~4.4 and Corollary~4.5]{FourthPaper2026}, with earlier Lean formalizations for abstract
functional DAGs in \cite[Lemma~3.3(b)]{MergeDepth2026}
(\leanname{MergeDepth/FieldDag.lean}, \leanname{MergeDepth.FunDag.cone}), for cone-size bounds on
skeletons and block machines of \leanname{Transpose.TM} in \cite[Lemma~4.7(a)]{MergeDepth2026}
(\leanname{MergeDepth/SkeletonCones.lean}, \leanname{MergeDepth.card_coneV_le} giving
$\sum_{h \le D}(k+1)^h$; and \leanname{MergeDepth/BlockMachine.lean},
\leanname{MergeDepth.blockMachine_card_cone_le}), and for native cut replay in \cite{AdaptiveTransfer2026}
(\leanname{AdaptiveTransfer/FrozenCutInvoice.lean},
\leanname{AdaptiveTransfer.stepValue_eq_of_uncut_native_cut},
\leanname{AdaptiveTransfer.native_cut_recovers}). What is new in \Cref{lem:tm-cut-cone} below are
the $(M.k + 1)^D$ cone-size instance directly on the native step DAG (\lean{TMDepStep}: control-state
spine plus last-work-visit edges) of frozen trace/cut fibres and the rooted $(1 + M.k)$-ary
symbolic tree representation (\lean{CausalConeTree} in \leanfile{MachineSweepConsequences.lean}),
extracted relative to column $c$ (\lean{traceConeTree}) and evaluated bottom-up on $\rot_c(X)$
directly through $M.\delta$ (\lean{evalConeTree}, syntactic support \lean{treeCellSupport}):

\begin{lemma}[Syntactic cut-cone support and symbolic tree evaluation on frozen fibres; cf.\ {\cite[Lemma~4.4 and Corollary~4.5]{FourthPaper2026}}, {\cite[Lemmas~3.3(b) and~4.7(a)]{MergeDepth2026}}, {\cite{AdaptiveTransfer2026}}]\label{lem:tm-cut-cone}
\formal{NativeSweepInterfaceProofs.lean}{tmTraceCutConeDetermination}
\formal{MachineSweepConsequences.lean}{evalConeTree_cellSupport,evalConeTree_entryCellSupport_card_le,evalConeTree_traceConeTree_eq}
Let $M$ be a \leanname{Transpose.TM} with $1 \le M.k$, let $d \ge 1$ and $D \ge 1$, and let
$C \subseteq \Sigma_d^{d\times d}$ satisfy \lean{IsFrozenTraceCutFibre} such that every $\Gamma$-free
causal path $\pi$ has at most $D$ vertices ($\pi.\mathrm{length} \le D$).
\begin{enumerate}[label=(\alph*)]
\item For every step $v < T$, $\abs{\lean{trueConeSteps}(\mathrm{workPos}, T, \Gamma, v)} \le
  (M.k + 1)^D$. For every output bit $(r, c, \ell)$ emitted at step
  $v = \mathrm{emitStep}(\lean{outputBitIdx}(d, r, c, \ell))$, the explicit syntactic cone support
  $Q_{\mathrm{bit}} := \lean{cutConeInputCells}(d, T, \mathrm{inPos}, \mathrm{workPos}, \Gamma, v)$
  consists of those cells visited by $\mathrm{inPos}$ on \lean{trueConeSteps}, satisfies
  $\abs{Q_{\mathrm{bit}}} \le (M.k + 1)^D$, and uniquely determines $X_{c,r,\ell}$ across $C$:
  whenever $X_1, X_2 \in C$ agree on $Q_{\mathrm{bit}}$, one has $X_1(c,r,\ell) = X_2(c,r,\ell)$.
\item Every depth-$D$ \lean{CausalConeTree} has $\abs{\lean{treeCellSupport}(\mathrm{tree})} \le
  (M.k + 1)^D$ and depends on $\rot_c(X)$ only through $\lean{treeCellSupport}(\mathrm{tree})$. Given
  any cut-value assignment $\mathrm{cutVal}$ agreeing with the frozen $\Gamma$-cut states and work
  symbols on $C$, evaluating $\lean{traceConeTree}$ at depth $D$ on $\rot_c(X)$ ($X \in C$) via
  \lean{evalConeTree} reproduces the post-step control state, written work symbols, and emitted output
  bit $X_{c,r,\ell}$ of the live machine run. Across all $b = \leanname{Transpose.entryBits}(d)$
  bit-lane trees of an output entry, the union of \lean{treeCellSupport} has size at most
  $\Delta = b(M.k + 1)^D$.
\end{enumerate}
\end{lemma}
\begin{outline}
Each step $w < T$ has at most $1 + M.k$ direct causal predecessors ($w - 1$ for the control state and
at most one last-visit step per work tape), so inducting on the depth budget $m \le D$ gives the
geometric sum bound $1 + (M.k+1)\bigl((M.k+1)^{m-1} - 1\bigr) < (M.k+1)^m$. Defining
$Q_{\mathrm{bit}} = \lean{cutConeInputCells}$ explicitly as the union of input cells read along
\lean{trueConeSteps} avoids self-referential singleton footprints: strong induction on $w < T$ shows
that inputs in $C$ agreeing on $Q_{\mathrm{bit}}$ induce identical transitions $M.\delta$ throughout
the $\Gamma$-stopped cone of $v$, and \lean{evalConeTree} replays those exact transitions bottom-up
from initial blank leaves and frozen $\Gamma$-cut leaves on $\rot_c(X)$.
\end{outline}

\subsection{Lean formalization of the \texorpdfstring{$1/(\Delta + 1)$}{1/(Delta + 1)} random-order bound for shared set codes}\label{sec:set-capacity}

For any recovery relation of maximum footprint size $\Delta$, the acyclic-subgraph bound of
Bar-Yossef, Birk, Jayram, and Kol \cite[Theorem~3]{BBJK2011} combined with a uniformly random vertex
ordering gives a maximum acyclic induced subgraph of size at least $(\text{\#clean})/(\Delta + 1)$
and hence rate at most $1 - 1/(\Delta + 1)$ up to self-referential entries. This elementary bound has
been in the series since \cite[Proposition~4.14(a)--(b)]{FourthPaper2026} (proved in text there, with
the $1 - 1/(\Delta+1)$ and acyclic-subgraph forms recalled in
\cite[Sections~4.1 and~6.7]{SeventeenthPaper2026}, and ordered-recovery injectivity formalized in
\leanname{LeanMachine.OrderedRecovery.eq_of_rank}). We record the Lean~4 formalization of its
column-class storage-code analogue in the notation of \Cref{def:bounded-arity-set-code}.

Given $p : \Z_d \to \mathrm{Fin}\,M$ and $Q : \mathrm{Fin}\,M \to \Z_d \to \mathrm{Finset}(\Z_d \times \Z_d)$,
evaluating $f_{p(c)}(\rot_c(X))_r$ reads the shifted footprint (\lean{shiftedFootprint})
\[
  F_{c,r} \;:=\; \bigl\{(r',\, c + u) \;:\; (r', u) \in Q_{p(c), r}\bigr\} \subseteq \Z_d \times \Z_d,
\]
and an input entry $(c, r)$ is \emph{clean} (\lean{cleanEntries}) if $(r, c) \notin \mathcal{E}$ and
$(c, r) \notin F_{c,r}$. Note that decoder sharing enters the following bound \emph{only} through the
$d M \Delta$ count of self-referential entries $(c, r) \in F_{c,r}$; once clean entries are isolated,
the random-order bound uses only the individual footprints $F_{c,r}$.

\begin{lemma}[Self-referential entries and clean-entry count]\label{lem:clean-count}
\formal{BoundedAritySetCapacityProofs.lean}{mem_shiftedFootprint_self_iff,selfReferentialAndCleanCount}
For all $d \ge 1$, $M, \Delta \in \N$, $p : \Z_d \to \mathrm{Fin}\,M$,
$Q : \mathrm{Fin}\,M \to \Z_d \to \mathrm{Finset}(\Z_d \times \Z_d)$, and
$\mathcal{E} \subseteq \Z_d \times \Z_d$:
\begin{enumerate}[label=(\alph*)]
\item $(c, r) \in F_{c,r}$ if and only if $(c, r - c) \in Q_{p(c), r}$.
\item If $\abs{Q_{m,r}} \le \Delta$ for all $(m, r)$, then $\abs{F_{c,r}} \le \Delta$ for every clean
  entry $(c, r)$ and
  \[
    \abs{\lean{cleanEntries}(p, Q, \mathcal{E})} \;\ge\; d^2 - \abs{\mathcal{E}} - d M \Delta.
  \]
\end{enumerate}
\end{lemma}
\begin{outline}
For each fixed $(m, r)$, the map $c \mapsto (c, r - c)$ is injective on $\Z_d$, so at most
$\abs{Q_{m,r}} \le \Delta$ columns $c \in p^{-1}(m)$ satisfy $(c, r) \in F_{c,r}$. Summing over
$m \in \mathrm{Fin}\,M$ and $r \in \Z_d$ bounds self-referential entries by $d M \Delta$.
\end{outline}

\begin{theorem}[Random-order capacity bound for bounded-arity shared set codes; after {\cite[Proposition~4.14(a)--(b)]{FourthPaper2026}} and {\cite[Theorem~3]{BBJK2011}}]\label{thm:random-order-capacity}
\formal{SetCapacityConsequences.lean}{orderedRecoverable_injective_and_card_le,exists_orderedRecoverableCells_bound,randomOrder_BoundedAritySetCapacity,corollary_OneClassArityOneHalf,corollary_OneClassArityTwoTwoThirds,asymptotic_SublinearClassesBrokenFractionRefutesRateOne}
Let $d \ge 1$, let $M, \Delta \in \N$, let $\Sigma$ be a finite alphabet with $\abs{\Sigma} \ge 2$,
and let $C \subseteq \Sigma^{d\times d}$ satisfy
$\lean{IsBoundedAritySetCode}(\Delta, p, \mathcal{E}, f, Q, C)$.
\begin{enumerate}[label=(\alph*)]
\item For every ranking $\mathrm{rk} : \Z_d \times \Z_d \to \N$, let
  $A(\mathrm{rk}) := \lean{orderedRecoverableCells}(p, Q, \mathcal{E}, \mathrm{rk})$ be the set of
  clean entries $(c, r)$ satisfying $\mathrm{rk}(y) < \mathrm{rk}(c, r)$ for all $y \in F_{c,r}$.
  Then restriction to $(\Z_d \times \Z_d) \setminus A(\mathrm{rk})$ is injective on $C$ (cf.\
  \leanname{LeanMachine.OrderedRecovery.eq_of_rank}), so
  $\abs{C} \le \abs{\Sigma}^{d^2 - \abs{A(\mathrm{rk})}}$, and averaging over all permutations of
  $\Z_d \times \Z_d$ produces an injective ranking $\mathrm{rk}_*$ with
  $(\Delta + 1)\abs{A(\mathrm{rk}_*)} \ge \abs{\lean{cleanEntries}(p, Q, \mathcal{E})}$.
\item Consequently, for every (empty or nonempty) code $C$,
  \[
    \frac{\log\abs{C}}{\log\abs{\Sigma}}
    \;\le\;
    d^2 \;-\; \frac{d^2 - \abs{\mathcal{E}} - d M \Delta}{\Delta + 1}.
  \]
  At $M = 1, \mathcal{E} = \emptyset$, this gives rate at most $\frac{1}{2} + \frac{1}{2d}$ at
  $\Delta = 1$ and $\frac{2}{3} + \frac{2}{3d}$ at $\Delta = 2$; and for any fixed $\Delta \ge 1$,
  $\delta \in (0, 1]$, any sequence of finite alphabets $(\Sigma_d)$ with $\abs{\Sigma_d} \ge 2$,
  $\abs{\mathcal{E}_d} \le (1 - \delta)d^2$, and $M_d = o(d/\Delta)$, the rate is eventually at most
  $1 - \frac{\delta}{2(\Delta + 1)}$.
\end{enumerate}
\end{theorem}
\begin{outline}
In~(a), strong induction on $\mathrm{rk}(c, r)$ shows that two codewords agreeing outside
$A(\mathrm{rk})$ agree on $A(\mathrm{rk})$ as well, because every $y \in F_{c,r}$ strictly precedes
$(c, r)$. For each clean entry $x = (c, r)$, $\abs{F_{c,r} \cup \{x\}} \le \Delta + 1$; for a
uniformly random permutation $\pi$ of $\Z_d \times \Z_d$, composing $\pi$ with the transposition
swapping $x$ and the $\pi$-maximum element of $F_{c,r} \cup \{x\}$ partitions the symmetric group so
that $x$ comes last in $F_{c,r} \cup \{x\}$ on at least a $1/(\Delta + 1)$ fraction of permutations,
and double counting selects $\mathrm{rk}_*$. Combining~(a) with \Cref{lem:clean-count}(b) and
$\log\abs{\Sigma} > 0$ gives~(b) (with $\log 0 = 0$ in Lean handling $C = \emptyset$).
\end{outline}

\begin{remark}[Comparison with small-window and unrestricted-arity bounds]\label{rem:capacity-context}
{\emergencystretch=4em
In \leanfile{BoundedAritySetCapacityProofs.lean} and \leanfile{SetCapacityConsequences.lean}
(\lean{oneHopRecoverableInjective}, \lean{existsRecoverableSubsetBound},
\lean{theorem1_BoundedAritySetCapacity}, \lean{growingClassesPositiveDeficit}, and
\lean{pMin_two_sided_theta_inv_arity}) we also formalized the weaker one-hop product-Bernoulli
variant with factor $\Delta^\Delta/(\Delta+1)^{\Delta+1} \in \bigl(1/(e(\Delta+1)),\,
1/(2(\Delta+1))\bigr]$ for $\Delta \ge 1$, together with explicit $M = 1, \mathcal{E} = \emptyset$ nonlinear lower bounds
for all $d \ge 2$ and $\abs{\Sigma} \ge 2$ (\lean{boundaryControl_DeficitVanishesAtDeltaEqD},
\lean{unrestrictedArity_RateLowerBoundAndTendstoOne}, and
\lean{unrestrictedArity_ExactLogarithmicAsymptoticAllDimensions}):
$\abs{C} = \lfloor\abs{\Sigma}^d / d\rfloor^d$ at cell arity $\Delta = d$ (where the clean-entry
coefficient $d^2 - dM\Delta$ vanishes), $\abs{C} \ge \abs{\Sigma}^{d(d-1)}$ at cell arity
$\Delta = 2d - 1$, and $\abs{C} \ge \abs{\Sigma}^{d^2}/d$ for prime $d$ and
$\abs{C} \ge \abs{\Sigma}^{d^2}/(2d)$ for all $d \ge 2$ at $\Delta = d^2$ (wrapping the full-view
orbit construction of \cite[Theorem~A.1 and Corollary~A.2]{SharedDecoders2026}, which reaches rate
$1 - o(1)$). At $\Delta \le 2$, stronger results (proved
in text, allowing a separate recovery function for each position) are already in
\cite{SeventeenthPaper2026}: at window $w = 1$, Proposition~6.21(b) of \cite{SeventeenthPaper2026}
gives the same exponent $\abs{C} \le q^{(d^2 + Md)/2}$ (rate $\frac{1}{2} + \frac{1}{2d}$ at
$M = 1$) without requiring shared functions, and at window $w \le 2$,
Proposition~6.22(b) of \cite{SeventeenthPaper2026} improves on the $\frac{2}{3} + \frac{2}{3d}$ bound
of \Cref{thm:random-order-capacity}(b) to $\abs{C} \le q^{d(d+1)/2 + d - 1}$ (rate
$\frac{1}{2} + O(1/d)$).
Meanwhile, Problem~9.4(b) of \cite{SeventeenthPaper2026} asks whether rate $\frac{1}{2} + o(1)$ holds
at fixed windows $w \ge 3$ over a \emph{fixed} alphabet of $q$ symbols, rather than merely rate
$< 1$.\par}
\end{remark}

\subsection{Negative brackets for two column-class cut reductions}\label{sec:sweep-remark}

\begin{remark}[Why subspace bounds and column-class set capacity stop short of Statement~(T)]\label{rem:sweep-negative}
It is natural to ask whether the linear shared-decoder bounds of
\Cref{sec:cubes,sec:sumsets,sec:linear-capacity,sec:aperiodic-six} or the set-capacity infrastructure
of \Cref{lem:tm-cut-cone} and \Cref{thm:random-order-capacity} can bound a periodic stretch of a
multitape Turing machine by conditioning on genie advice or a bounded-depth trace cut $\Gamma$. Three
distinct barriers prevent a direct reduction:
\begin{enumerate}[label=(\arabic*)]
\item \emph{Nonlinear set fibres versus linear subspaces
  (\cite[Problem~9.4(b)]{SeventeenthPaper2026}; \cite[Theorem~A.1 and Corollary~A.2]{SharedDecoders2026}).}
  Freezing genie advice or a cut transcript on a finite-alphabet machine selects a \emph{nonlinear
  subset} $C \subseteq \Sigma_d^{d\times d}$ rather than a linear subspace, so the linear decoder and
  Fourier-filtration bounds of \Cref{sec:cubes,sec:sumsets,sec:linear-capacity,sec:aperiodic-six} do
  not apply directly. Moreover, Theorem~A.1 and Corollary~A.2 of \cite{SharedDecoders2026} prove that
  shared set decoders with unbounded window reach rate $1 - o(1)$ ($\abs{C} \ge \abs{\Sigma}^{d^2}/(2d)$),
  so any set-level bridge must enforce a bounded read window.
\item \emph{Depth-$D$ cut cost exceeds the $1/(\Delta+1)$ deficit
  (\cite[Propositions~4.13 and~4.14(c)]{FourthPaper2026}, proved in text).}
  In \cite[Propositions~4.13 and~4.14(c)]{FourthPaper2026}, path depth counts edges (stated for
  oblivious machines, with the same spine-window counting argument applying to the common trajectory of
  a frozen-trace fibre), so a $\Gamma$-free path of at most $D$ vertices (\Cref{lem:tm-cut-cone}) has
  edge-depth at most $D - 1$: cutting all causal paths down to $D$ vertices over a run of length $T \ge L$
  at cut bit-cost $\eps L = \Omega(L/D)$ forces $(M.k + 1)^D > 1/\eps$ (evaluating
  \cite[Proposition~4.14(c)]{FourthPaper2026} at edge-depth $D - 1$), while the entry-level cone arity
  $\Delta = b(M.k + 1)^D$ grows exponentially in $D$ and includes $b = \max(1, \clog{d})$. Thus an
  $O(1/\Delta)$ rate deficit is always smaller than the cut-advice rate and, as proved in
  \cite[Proposition~4.14(c)]{FourthPaper2026}, cannot by itself contradict a transposer.
\item \emph{Per-entry cut labels along spine edges, and separated high-rate storage codes
  (\cite[commentary on Problem~9.4(b)]{SeventeenthPaper2026}; \cite{StorageCodes2026,HighRate2026}).}
  \Cref{prop:negative-brackets} below illustrates that two naive column-class reductions of a
  depth-$D$ $\Gamma$-cut fibre---discarding all cut-touching outputs, or positing a worst-case
  full-product class count $M_{\mathrm{cls}} := M_0\bigl((M.q+1)(M.g+1)^{M.k}\bigr)^{L_{\mathrm{cut}}}$
  across $L_{\mathrm{cut}}$ per-column cut slots---make the numerator $d^2 - \abs{\mathcal{E}} - d M \Delta$
  in \Cref{thm:random-order-capacity}(b) strictly negative. (Note that \Cref{prop:negative-brackets}(b)
  uses only $M_{\mathrm{cls}} \ge L_{\mathrm{cut}} + 1$ and does not bound how many cut patterns
  actually occur: on periodic stretches few cut patterns may occur, where the bracket need not be
  negative and the obstruction is instead item~(2) above.) In \Cref{prop:negative-brackets}, this
  occurs because \lean{TMDepStep} includes the control-state spine edge $u \to u + 1$, attaching cut
  labels to individual output steps rather than whole columns; this reinforces the observation in the
  commentary on \cite[Problem~9.4(b)]{SeventeenthPaper2026} that cut patterns divide outputs into
  boundedly many per-entry classes, motivating the per-position fixed-alphabet window formulation of
  \cite[Problem~9.4(b)]{SeventeenthPaper2026}. Meanwhile, linear-time transposers must use curved,
  head-dependent, non-piecewise-affine motion on interacting work tapes rather than uniform
  home-based passes (\cite[Theorems~3.10 and~3.12 and Corollary~3.24]{FourthPaper2026}), and
  \cite{StorageCodes2026,HighRate2026}
  construct bounded-degree paired-distance-separated storage codes above rate $\frac{1}{2}$ and
  approaching rate $1$, so graph degree and separation alone cannot rule out linear-time transposition
  without exploiting one-dimensional head trajectory geometry.
\end{enumerate}
\end{remark}

\begin{proposition}[Negative brackets for two specific column-class cut reductions]\label{prop:negative-brackets}
\formal{MachineSweepConsequences.lean}{perColumnCutAbsorption_bracket_neg}
Let $M$ be a \leanname{Transpose.TM} with $1 \le M.k$, let $d \ge 1$ and $D \in \N$, and let
$C \subseteq \Sigma_d^{d\times d}$ be a \emph{nonempty} fibre satisfying \lean{IsFrozenTraceCutFibre}
whose $\Gamma$-free causal paths have at most $D$ vertices. Set
$\Delta := \leanname{Transpose.entryBits}(d) \cdot (M.k + 1)^D$.
\begin{enumerate}[label=(\alph*)]
\item \emph{Discarding cut-touching outputs}: let $\mathcal{E}_{\mathrm{cut}}$
  (\lean{cutTouchingOutputCells} for $d$, $T$, $\mathrm{workPos}$, $\mathrm{emitStep}$, and $\Gamma$)
  be the set of output coordinates $(r, c)$ having a bit cone that meets $\Gamma$ along a direct
  \lean{TMDepStep}. Then $d^2 \le \abs{\mathcal{E}_{\mathrm{cut}}} + D$, and for every $M_0 \ge 1$
  and every $\mathcal{E}_{\mathrm{sched}} \subseteq \Z_d \times \Z_d$, setting
  $\mathcal{E} := \mathcal{E}_{\mathrm{sched}} \cup \mathcal{E}_{\mathrm{cut}}$ gives
  \[
    d^2 - \abs{\mathcal{E}} - d M_0 \Delta \;<\; 0.
  \]
\item \emph{Worst-case full-product class refinement across per-column cut slots}: suppose every cut
  step $u \in \Gamma$ touched by a bit cone of an unbroken output $(r, c) \notin \mathcal{E}_{\mathrm{sched}}$
  lies among $L_{\mathrm{cut}}$ column-indexed cut slots $\mathrm{colCutStep}(c) : \mathrm{Fin}\,L_{\mathrm{cut}} \to \N$.
  Then $d^2 \le \abs{\mathcal{E}_{\mathrm{sched}}} + D + d L_{\mathrm{cut}} D$, and for every
  $M_0 \ge 1$, whenever $(M.q + 1)(M.g + 1)^{M.k} \ge 2$, the worst-case full-product class count
  $M_{\mathrm{cls}} := M_0 \bigl((M.q + 1)(M.g + 1)^{M.k}\bigr)^{L_{\mathrm{cut}}}$ satisfies
  \[
    d^2 - \abs{\mathcal{E}_{\mathrm{sched}}} - d M_{\mathrm{cls}} \Delta \;<\; 0.
  \]
\end{enumerate}
\end{proposition}
\begin{outline}
Since $C \ne \emptyset$, the $d^2$ lane-$0$ emission steps $v(r,c) = \mathrm{emitStep}(\mathrm{idx}_0(r,c))$
(via \lean{outputBitIdx} at lane~$0$)
are pairwise distinct. For any $(r,c)$ with $v(r,c) \ge D$, the descending control-state spine from
$v(r,c)$ cannot avoid $\Gamma$ for $D$ steps without creating a $\Gamma$-free path of $D + 1$
vertices; taking the largest $u \in \Gamma \cap [v(r,c) - D, v(r,c))$ places $u + 1$ in the
$\Gamma$-cut cone of $v(r,c)$ with $u \to u + 1$, so $(r,c) \in \mathcal{E}_{\mathrm{cut}}$. At most
$D$ emission steps satisfy $v(r,c) < D$, giving $d^2 - \abs{\mathcal{E}_{\mathrm{cut}}} \le D <
2^D \le \Delta \le d M_0 \Delta$ and proving~(a). In~(b), each slot $u = \mathrm{colCutStep}(c, s)$
can serve as that most recent cut step for at most $D$ emission steps $v(r,c) \in (u, u + D]$, leaving
at most $D + d L_{\mathrm{cut}} D \le d(L_{\mathrm{cut}} + 1)D$ unbroken entries; meanwhile
$(M.q + 1)(M.g + 1)^{M.k} \ge 2$ implies $M_{\mathrm{cls}} \ge 2^{L_{\mathrm{cut}}} \ge L_{\mathrm{cut}} + 1$,
and $D < \Delta$ yields $d(L_{\mathrm{cut}} + 1)D < d M_{\mathrm{cls}} \Delta$.
\end{outline}

\subsection{Open problems}\label{sec:open}

Statement~(T) remains open. Within the models studied in this paper and in
\cite{SeventeenthPaper2026,ClassUniformMemo2026,SharedDecoders2026}, four concrete algebraic and
combinatorial problems remain unresolved:
\begin{enumerate}[label=\arabic*.]
\item \textbf{Macroscopic rank deficit for $M_1 \ge 2$ and positive characteristic in two-stage linear composition (\cite[Sections~6 and~8]{SharedDecoders2026}).}
  For fixed $M_1 \ge 2$ and $M_2 \ge 1$ (at $M_1 = 1$, \cite[Theorem~6.6]{SharedDecoders2026} already
  gives $\Omega(bd^2 / M_2)$ for $d \ge 4M_2 + 1$), can the $\Omega\bigl(bd(\log d)^{1/(2M_2 - 1)}\bigr)$ error-rank bound
  of \Cref{thm:2b} be strengthened to a positive normalized deficit $\Omega_{M_1,M_2}(bd^2)$ as asked
  in \cite[Section~8]{SharedDecoders2026}, and extended to fields in which $\operatorname{char}(F) \mid d$
  (such as $\F_2$ at $d = 2^m$)?
\item \textbf{Closing $\frac{19}{32}$ to $\frac{1}{2} + o(1)$ and positive characteristics in Problem~9.1(a)--(b) of \cite{SeventeenthPaper2026}.}
  Can the $\frac{19}{32}\,bd^2 + \frac{3}{2}\,bd^{3/2}$ two-class ceiling of \Cref{thm:1b-two-class}
  be lowered to $(\frac{1}{2} + o(1))bd^2$, can the averaged multi-class bound of
  \Cref{thm:1b-multi-class} ($M \ge 3$) be extended from primes to composite $d$, and can
  \Cref{thm:1a-single-shift,thm:1b-two-class,thm:1b-multi-class} be extended to fields with
  $\operatorname{char}(F) \mid d$ (\cite[Problem~9.1(a)]{SeventeenthPaper2026})? And for
  $\{0,1,3\} \mid \{2,4,5\}$ on $\Z_6$ (\Cref{thm:aperiodic-six-25}), what is
  $D_1(\{0,1,3\},\{2,4,5\})$ over fields of characteristic~$2$ ($25 \le D_1 \le 28$) and
  characteristic~$3$ ($21 \le D_1 \le 28$ from symmetric matrices and
  \cite[Theorem~6.4]{SeventeenthPaper2026}; whether $25$ is attainable in characteristic~$3$ is
  \cite[Problem~6.1]{ClassUniformMemo2026})?
\item \textbf{Remaining dimensions in Remark~5.6 and two-sided multi-scale bounds in Problem~9.3 of \cite{SeventeenthPaper2026}, and periodic Walsh masks.}
  Does any three-term swap code exceed $d(d+1)/2$ at $d = 5$ or at any $d \ge 8$ outside $\{9, 12\}$
  (\cite[Remark~5.6]{SeventeenthPaper2026})? Is $\bigl\lfloor \frac{1}{2}\bigl(d^2 +
  \sum_{i, j} \min(\abs{I_i}, \abs{J_j})\bigr)\bigr\rfloor$ (\Cref{rem:two-sided-numerical}) an upper
  bound on swap-code dimension for all two-sided interval partitions, and can
  \Cref{thm:two-sided-coarse-fine} be extended from two scales to general multi-scale dyadic interval
  profiles (\cite[Problem~9.3]{SeventeenthPaper2026})? On the periodic variant of
  \cite[Problem~9.3]{SeventeenthPaper2026} ($d = 2^m$), when both $A_a$ and $B_a$ lie in the at most
  $(m + 1)2^m$-dimensional span of periodic Walsh digit masks composed with rotations
  (\Cref{prop:walsh-dichotomy}), is $\dim_F {}$\leanname{SwapRigidity.swapCode}$(A, B) \le (1 - c)d^2$
  for a constant $c > 0$?
\item \textbf{Fixed-alphabet windows, oblivious/adaptive interfaces, and multi-scale invariants (\cite[Problems~9.4 and~9.5]{SeventeenthPaper2026}).}
  Over a fixed finite alphabet of $q$ symbols and a fixed window $w \ge 3$ on a spread one-class
  relation, does the $\frac{1}{2} + o(1)$ rate bound hold with a separate recovery function for each
  position or a single shared function (\cite[Problem~9.4(b)]{SeventeenthPaper2026}), or in its
  whole-row $\rho$-code form (\cite[Problem~9.4(c)]{SeventeenthPaper2026})? On the machine side,
  \cite[Problem~9.4(a)]{SeventeenthPaper2026} asks for a cone-system analogue of Theorem~TP on
  \emph{oblivious} schedules (while \cite[Theorem~H.5 and Corollary~H.6]{AdaptiveTransfer2026} transfer
  entry-relation storage bounds to cut runs of \emph{adaptive} transposers), and
  \cite[Problem~9.5]{SeventeenthPaper2026} asks for a quantity subadditive over machine steps, large for
  the swap at every scale separately, and \emph{``not raised at all scales by a single linear-time
  operation''}.
\end{enumerate}
